\documentclass[11pt]{article}

% write-latex compliant: no fontspec, no XeLaTeX-only packages.
\usepackage[utf8]{inputenc}
\usepackage[T1]{fontenc}
\usepackage{lmodern}
\usepackage[margin=1in]{geometry}
\usepackage{amsmath, amssymb, amsthm, mathtools}
\usepackage{hyperref}
\hypersetup{colorlinks=true, linkcolor=blue, urlcolor=blue, citecolor=blue}

\newtheorem{theorem}{Theorem}
\newtheorem{lemma}[theorem]{Lemma}
\newtheorem{corollary}[theorem]{Corollary}
\theoremstyle{definition}
\newtheorem{definition}[theorem]{Definition}
\theoremstyle{remark}
\newtheorem{remark}[theorem]{Remark}

\title{ {{improvement_title}} }
\author{Companion proof for: \href{ {{repo_url}} }{ {{github_username}}/{{slug}} } \\ \small Improvement slug: \texttt{ {{improvement_slug}} }}
\date{}

\begin{document}
\maketitle

\iffalse
NOTE: this file is a per-improvement proof, generated by the study-paper skill, Stage 7.
PROOF mode is selected when the improvement is mathematical, theoretical, or a
complexity-class statement (see Stage 7 heuristic). The proof must be complete:
no "obvious", no "left as exercise", no "by standard arguments".
Mustache placeholders {{...}} are replaced per-improvement before compilation.
\fi

\begin{abstract}
{{abstract_paragraph}}
\end{abstract}

\section*{Setup and notation}

{{setup_paragraph}}

\begin{definition}[Working notation]
{{notation_block}}
\end{definition}

\section*{Theorem}

\begin{theorem}[ {{theorem_name}} ]
\label{thm:main}
{{theorem_statement}}
\end{theorem}

\section*{Proof}

\begin{proof}
{{proof_body}}
\end{proof}

\section*{Discussion}

\paragraph{What this gives us.}
{{discussion_what_it_gives}}

\paragraph{What this does not give us.}
{{discussion_what_it_does_not}}

\paragraph{Where it would break.}
{{discussion_failure_modes}}

\paragraph{Empirical companion.}
This proof is the validation file for the improvement \texttt{ {{improvement_slug}} } in \texttt{05-improvements.tex}. If a measurement-mode prototype also exists at \texttt{improvements/{{improvement_slug}}.py}, it serves as a sanity check, not as evidence --- the theorem above is the load-bearing claim.

\end{document}
